\documentclass[10pt]{article}
\usepackage[utf8]{inputenc}
\usepackage[T1]{fontenc}
\usepackage{amsmath}
\usepackage{amsfonts}
\usepackage{amssymb}
\usepackage[version=4]{mhchem}
\usepackage{stmaryrd}
\usepackage{bbold}

\begin{document}
$29$

\section*{L9: Arithmetic Functions and}
\section*{Multiplicativity}
\begin{itemize}
  \item A comment about notation: $\sum f(d)$ means the $d(n, d>0$\\
summation of $f$ over all distinct positive divisors $d$ of $n$. For example,
\end{itemize}

$$
\sum_{\substack{d \mid 12 \\ d>0}} f(d)=f(1)+f(2)+f(3)+f(4)+f(6)+f(12) .
$$

Theorem 3.1: Let $f$ be an arithmetic function, and for $n \in \mathbb{Z}$ with $n>0$, let

$$
F(n)=\sum_{d \mid n, d>0} f(d)
$$

If $f$ is multiplicative, then $F$ is multiplicative.\\
Proof: Let $m$ and $n$ be relatively prime positive integers. To prove that $F$ is multiplicative, we must prove that $F(m n)=F(m) F(n)$. We have that

$$
F(m n)=\sum_{d \text { m } n, d>0} f(d) \text {. Claim: Since }(m, n)=1,
$$

each divisor $d>0$ of $m n$ can be written uniquely as $d_{1} d_{2}$ where $d_{1}, d_{2}>0$, $d_{1}\left|m, d_{2}\right| n$, and $\left(d_{1}, d_{2}\right)=1$, and each such product $d_{1} d_{2}$ corresponds to a divisor of $m n$. The case when either $m=1$ or $n=1$ are trivial. Assume that $m>1$ and $n>1$, and let $m=p_{1}^{a_{1}} p_{2}^{a_{2}} \cdots p_{r}^{a_{r}}$

$$
n=q_{1}^{b_{1}} q_{2}^{b_{2}} \cdots q_{s}^{b_{s}}
$$

be the prime factorisations of $m$ and $n$ respectively. Since $(m, n)=1$, all $p_{i}$ and $q_{j}$ are distinct. Let $d$ be a positive divisor of $m n$. Then

$$
d=p_{1}^{e_{1}} p_{2}^{e_{2}} \cdots p_{r}^{e_{r}} q_{1}^{f_{1}} q_{2}^{f_{2}} \cdots q_{s}^{f_{s}},
$$

where the integers $e_{i}$ and $f_{j}$ satisfy\\
$0 \leqslant e_{i} \leqslant a_{i}$ and $0 \leqslant f_{j} \leqslant b_{j}$ for all $i$ and $j$. Now $d_{1}=p_{1}^{e_{1}} p_{2}^{e_{2}} \cdots p_{r}^{e_{r}}$ and $d_{2}=q_{1}^{f_{1}} q_{2}^{f_{2}} \ldots q_{s}^{f_{s}}$ are as desired. The fact\\
that each such product corresponds to a divisor of $d$ is clear. This proves the claim. Then,

$$
\begin{aligned}
F(m n) & =\sum_{\substack{d / m n \\
d>0}} f(d) \\
& =\sum_{\substack{d_{1}\left|m, d_{2}\right| n \\
d_{1}, d_{2}>0}} f\left(d_{1} d_{2}\right) \\
& =\sum_{\substack{d_{1}\left|m, d_{2}\right| n \\
d_{1}, d_{2}>0\\
}} f\left(d_{1}\right) f\left(d_{2}\right) \quad\binom{\text { since } f \text { is }}{\text { multiplicative }} \\
& =\left(\sum_{d_{1} \mid m} f\left(d_{1}\right)\right) \cdot\left(\sum_{d_{2} \mid n} f\left(d_{2}\right)\right) \\
& =F(m) F(n),
\end{aligned}
$$

as desived.

Example: Let $m=3$ and $n=4$, and keep the notation the same as Theorem 3.1. We show\\
that $F(3.4)=F(3) F(4)$. We have

$$
\begin{aligned}
F(3.4)= & \sum_{d \mid 12, d>0} f(d) \\
= & f(1)+f(2)+f(3)+f(4)+f(6)+f(12) \\
= & f(1)+f(2)+f(4)+f(3)+f(6)+f(12) \\
= & f(1.1)+f(1.2)+f(1.4)+f(3.1) \\
& +f(3.2)+f(3.4) \\
= & f(1) f(1)+f(1) f(2)+f(1) f(4) \\
& +f(3) f(1)+f(3) f(2)+f(3) f(4) \\
= & (f(1)+f(3))(f(1)+f(2)+f(4)) \\
= & \left(\sum_{d_{1} \mid 3, d_{1}>0} f\left(d_{1}\right)\right) \cdot\left(\sum_{d_{2} \mid 4, d_{2}>0} f\left(d_{2}\right)\right) \\
= & F(3) F(4),
\end{aligned}
$$

as desived.

$$
\text { The Euler } \phi \text {-function }
$$

Theorem 3.2: The Euler $\phi$-function is multiplicative.

Proof: Let $m$ and $n$ be relatively prime positive integers. We must prove that $\phi(m n)=\phi(m) \phi(n)$. We display the positive integers not exceeding mn as a matrix:

$$
\begin{array}{ccccc}
1 & m+1 & 2 m+1 & \cdots & (n-1) m+1 \\
2 & m+2 & 2 m+2 & \cdots & (n-1) m+2 \\
\vdots & \vdots & \vdots & \vdots & \\
i & m+i & 2 m+i & \cdots & (n-1) m+i \\
\vdots & \vdots & \vdots & \vdots & \vdots \\
m & 2 m & 3 m & \cdots & n m
\end{array}
$$

Consider the ith row of the matrix above. If $(m, i)>1$, then no entry of the ith row is relatively prime to $m$, and so not velatively. prime to $m n$ (every entry of the ith row is of the form $k m+i$ where $k \in \mathbb{C}$, so if $(m, i)=d>1$, then $d / m$ and $d / i$ implies that d lem + i by Proposition 1.2). Hence, to determine how many integers in the matrix that are\\
relatively prime to $m n$, which is $\phi(m n)$, it suffices to examine the ith row of the matrix if and only if $(m, i)=1$. There are $\phi(m)$ such rows. Now, given such an ith row, we must determine how many integers are relatively prime to $m n$. The entries of such a row are

$$
i, m+i, 2 m+i, \cdots,(n-1) m+i
$$

Since $(m, i)=1$, each of these integers are relatively prime to $m$. Furthermore, these $n$ integers form a complete residue system modulon.\\
Otherwise, $j m+i \equiv k m+i(\bmod n)$ for some\\
$0 \leq j, k \leq n-1$ with $j \neq k$. Then $j m \equiv k m(\bmod n)$, and since $(m, n)=1$, the inverse of $m$ modulon exists, say $m^{\prime}$, and we have jmm' $\equiv k m m^{\prime}(\bmod n)$, and so $j \equiv k(\bmod n)$. Impossible! Thus exactly $\phi(n)$ of these integers are relatively prime to $n$. Since these $\phi(n)$ integers are relatively prime to $m$, they are relatively prime to $m n$. In summary, exactly $\phi(m)$ rows of the matrix\\
contain integers relatively prime to m. Each (7) such row contains exactly $\phi(n)$ integers relatively prime to $m n$. Thus $\phi(m n)=\phi(m) \phi(n)$, as required.

\begin{itemize}
  \item Since $\phi$ is multiplicative, it is determined by its values at powers of prime numbers.
\end{itemize}

Theorem 3.3: Let $p$ be a prime number and let $a \in \mathbb{Z}$ with $a>0$. Then

$$
\phi\left(p^{a}\right)=p^{a}-p^{a-1} .
$$

Proof: The total number of positive integers not exceeding $p^{a}$ is $p^{a}$. The positive integers not exceeding $p^{a}$ that are not relatively prime to $p^{a}$ are precisely the multiples of $P$ :

$$
p, 2 p, 3 p, \cdots,\left(p^{a-1}\right) p .
$$

There are $p^{a-1}$ such multiples, and so we conclude that $\phi\left(p^{a}\right)=p^{a}-p^{a-1}$.\\
Theovem 3.4: Let $n \in \mathbb{Z}$ with $n>0$. Then

$$
\phi(n)=n \prod_{p \mid n, p \text { prime }}\left(1-\frac{1}{p}\right) .
$$

\begin{itemize}
  \item The notation TT means the product is (8) pln,p prime\\
taken over all distinct prime divisors $p$ of $n$.\\
For example,
\end{itemize}

$$
\pi\left(1-\frac{1}{p}\right)=\left(1-\frac{1}{2}\right)\left(1-\frac{1}{3}\right)
$$

pli2, p prime\\
Proof (of Theorem 3.4): The result is clear for $n=1$ (with the convention that the empty product is taken to be 1). Assume that $n>1$ and that $n=p_{1}^{a_{1}} p_{2}^{a_{2}} \ldots p_{r}^{a_{r}}$ with $p_{1}, p_{2}, \ldots, p_{r}$ dirtinct prime numbers and $a_{1}, \ldots, a_{r}$ positive integers. Then

$$
\begin{aligned}
\phi(n) & =\phi\left(p_{1}^{a_{1}} p_{2}^{a_{2}} \cdots p_{r}^{a_{r}}\right), \text { by Theorem } 3.2 \\
& =\phi\left(p_{1}^{a_{1}}\right) \phi\left(p_{2}^{a_{2}}\right) \cdots \phi\left(p_{r}^{a_{r}}\right) \text { by Theorem } 3.3 \\
& =\left(p_{1}^{a_{1}}-p_{1}^{a_{1}-1}\right)\left(p_{2}^{a_{2}}-p_{2}^{a_{2}-1}\right) \cdots\left(p_{r}^{a_{r}}-p_{r}^{a_{r}-1}\right) \\
& =p_{1}^{a_{1}}\left(1-\frac{1}{p_{1}}\right) p_{2}^{a_{2}}\left(1-\frac{1}{p_{2}}\right) \cdots p_{r}^{a_{r}}\left(1-\frac{1}{p_{r}}\right) \\
& =p_{1}^{a_{1}} p_{2}^{a_{2}} \cdots p_{r}^{a_{r}}\left(1-\frac{1}{p_{1}}\right)\left(1-\frac{1}{p_{2}}\right) \cdots\left(1-\frac{1}{p_{r}}\right) \\
& =n \prod_{1 n}\left(1-\frac{1}{p}\right) .
\end{aligned}
$$

Example: Compute $\phi(504)$. We have $504=2^{3} 3^{2} 7$,\\
so $\phi(504)=504\left(1-\frac{1}{2}\right)\left(1-\frac{1}{3}\right)\left(1-\frac{1}{7}\right)$

$$
\begin{aligned}
& =504(1 / 2)(2 / 3)(6 / 7) \\
& =144 .
\end{aligned}
$$

Theorem 3.5 (Gauss): Let $n \in \mathbb{Z}$ with $n>0$. Then\\
$\sum \phi(d)=n$.\\
$d \mid n, d>0$\\
Proof: Let $d$ be a positive divisor of $n$ and put

$$
S_{d}:=\{m \in \mathbb{Z}: 1 \leq m \leq n,(m, n)=d\} .
$$

Now $(m, n)=d$ if and only if $\left(\frac{m}{d}, \frac{n}{d}\right)=1$. Thus the number of integers in $S_{d}$ is equal to the number of positive integers not exceeding $n / d$ that are relatively prime to $n / d$ i.e.\\
$\left|S_{d}\right|=\phi(n / d)$. Since every integer from 1 to $n$ inclusive belongs to exactly one $S_{d}$, we have that $n=\sum \phi(n / d)$. As $d$ runs through $d \mid n, d>0$\\
all positive divisors of $n$, so to does $n / d$, thus

$$
n=\sum_{d \mid n, d>0} \phi(n / d)=\sum_{d \ln , d>0} \phi(d),
$$

as desived.\\
Example: Let $n=12$ and all the notation be as in Theorem 3.5. We must show that

$$
\sum_{d \mid 12, d>0} \phi(d)=12 .
$$

We have

$$
\begin{array}{ll}
S_{1}=\left\{1, S_{1}, 11\right\}, & \left|S_{1}\right|=4=\phi(12) \\
S_{2}=\{2,10\}, & \left|S_{2}\right|=2=\phi(6) \\
S_{3}=\{3,9\}, & \left|S_{3}\right|=2=\phi(4) \\
S_{4}=\{4,8\}, & \left|S_{4}\right|=2=\phi(3) \\
S_{6}=\{6\}, & \left|S_{6}\right|=1=\phi(2) \\
S_{12}=\{12\} . & \left|S_{12}\right|=1=\phi(1)
\end{array}
$$

In all cases we have that $\left|S_{d}\right|=\phi(12 / d)$. Since each integer from 1 to 12 inclusive\\
belongs to exactly one $S_{d}$ we have that (II)

$$
12=\phi(1)+\phi(2)+\phi(3)+\phi(4)+\phi(6)+\phi(12) \text {. }
$$


\end{document}